```latex
\documentclass[11pt]{amsart}

\usepackage{amsmath,amssymb,amsthm}
\usepackage[colorlinks=true,citecolor=blue,linkcolor=blue,urlcolor=blue]{hyperref}

\newcommand{\Sph}{\mathbb S}
\newcommand{\Chat}{\widehat{\mathbb C}}
\newcommand{\Teich}{\operatorname{Teich}}
\newcommand{\Mod}{\operatorname{Mod}}
\newcommand{\MCG}{\operatorname{MCG}}
\newcommand{\Def}{\operatorname{Def}}
\newcommand{\Per}{\operatorname{Per}}

\theoremstyle{plain}
\newtheorem{theorem}{Theorem}
\newtheorem{lemma}{Lemma}

\theoremstyle{remark}
\newtheorem{remark}{Remark}

\begin{document}

\title{A Dehn twist with an infinite orbit on components of the period-four deformation space}
\author{}
\date{}
\maketitle

\begin{abstract}
We identify an explicit curve for the deformation space associated to a quadratic
rational map with one critical point of period four.  In the notation
\(p_j=f^{\circ j}(c_1)\), \(j=0,1,2,3\), and \(v=f(c_2)\), the curve is any
essential simple closed curve separating \(\{p_1,v\}\) from
\(\{p_0,p_2,p_3\}\).  A suitable positive power of the Dehn twist about this
curve preserves the deformation space and has an infinite orbit on its connected
components.
\end{abstract}

\section*{Statement and proof}

Let
\[
p_0=c_1,\qquad p_1=f(c_1),\qquad p_2=f^{\circ 2}(c_1),\qquad
p_3=f^{\circ 3}(c_1),
\]
so that \(p_0\mapsto p_1\mapsto p_2\mapsto p_3\mapsto p_0\), and let
\(v=f(c_2)\notin \{p_0,p_1,p_2,p_3\}\).  Thus
\[
A=\{p_0,p_1,p_2,p_3\},\qquad B=A\cup\{v\}.
\]

\begin{theorem}
Let \(\gamma\subset \Sph^2\setminus B\) be a simple closed curve separating
\[
\{p_1,v\}\quad\text{from}\quad \{p_0,p_2,p_3\}.
\]
Let \(T_\gamma\) be the positive Dehn twist about \(\gamma\).  Then a positive
power of \(T_\gamma\), in fact \(T_\gamma^2\), preserves the deformation space
\(\Def(f,A,B)\).  Moreover the action of \(T_\gamma^2\) on the set of connected
components of \(\Def(f,A,B)\) has an infinite orbit.
\end{theorem}

\begin{proof}
The assertion is invariant under Möbius conjugacy, so we use the standard
normal form
\[
p_0=0,\qquad p_1=\infty,\qquad p_2=1,\qquad p_3=\rho.
\]
Then
\[
f=f_\rho(z)=1+\frac{b}{z}+\frac{c}{z^2},
\]
with
\[
b=\frac{-\rho^2-\rho+1}{\rho-1},\qquad
c=\frac{2\rho^2-\rho}{\rho-1},
\]
and
\[
v=v_\rho
=1-\frac{(\rho^2+\rho-1)^2}{4(\rho-1)(2\rho^2-\rho)}.
\]
The curve \(\gamma\) is the curve enclosing precisely the two marked points
\(\{\infty,v\}\).

First, \(\gamma\) is an equalizing curve in the sense of Firsova--Kahn--Selinger:
after forgetting \(v\), the curve becomes peripheral, and every non-peripheral
component of \(f^{-1}(\gamma)\) disappears after forgetting
\(f^{-1}(B)\setminus A\).  Equivalently,
\[
i_*(\gamma)=\varnothing,\qquad f^*(\gamma)=\varnothing
\]
as multicurves on \((\Sph^2,A)\).  This is exactly case (4) in the classification
of equalizing curves for \(\Per'_4\) in Firsova--Kahn--Selinger
\cite[Appendix, case (4)]{FKS}.  Since \(\deg f=2\), the square \(T_\gamma^2\)
is liftable; its lift and its forgetful image are both trivial in
\(\MCG(A)\).  Hence \(T_\gamma^2\) is a special liftable mapping class, and so
it preserves \(\Def(f,A,B)\).

It remains to prove that its action on components has an infinite orbit.  Let
\[
\eta
\]
be the simple closed curve separating \(\{0,1\}\) from
\(\{\infty,\rho,v\}\).  The curves \(\gamma\) and \(\eta\) are disjoint.  Thus
their Dehn twists commute.

Let \(D_0\) denote the connected component of \(\Def(f,A,B)\) containing the
base complex structure, and let \(E_f\) be the stabilizer of \(D_0\) inside the
group of special liftables.  Hironaka--Koch prove that the multitwist
\[
M:=T_\gamma T_\eta
\]
has no non-zero power in \(E_f\):
\[
M^n\notin E_f\qquad\text{for every }n\neq 0.
\]
In their notation this is the element constructed in Section 5.3; it is shown
to be special liftable in \cite[Proposition 5.6]{HK} and to have no non-zero
power in \(E_f\) in \cite[Proposition 5.8]{HK}.  The identification of their
loop with \(T_\gamma T_\eta\) is the standard Picard--Lefschetz monodromy
identification: a positive meridian about the boundary divisor of
\(\overline{M}_{0,5}\) determined by a separating curve is the corresponding
positive Dehn twist.

Also, \(T_\eta\) stabilizes \(D_0\).  Indeed, the moduli image of the deformation
space is \(\Per'_4\), parametrized by \(\rho\), and the puncture \(\rho=\infty\)
corresponds precisely to pinching \(\eta\).  A small positive loop about this
puncture gives the Dehn twist \(T_\eta\), so \(T_\eta\in E_f\).

Now suppose, for contradiction, that the orbit of \(D_0\) under \(T_\gamma^2\)
were finite.  Then for some \(n>0\),
\[
T_\gamma^{2n}\in E_f.
\]
Since \(T_\eta\in E_f\) and \(T_\gamma\) commutes with \(T_\eta\), we would have
\[
M^{2n}=(T_\gamma T_\eta)^{2n}
      =T_\gamma^{2n}T_\eta^{2n}\in E_f,
\]
contradicting the Hironaka--Koch theorem above.  Therefore no non-zero power of
\(T_\gamma^2\) stabilizes \(D_0\), and the orbit
\[
\{T_\gamma^{2n}D_0:n\in\mathbb Z\}
\]
is infinite.
\end{proof}

\begin{remark}
In the original, unnormalized notation, the curve is simply:
\[
\boxed{\text{a simple closed curve separating } \{f(c_1),\,f(c_2)\}
\text{ from } \{c_1,\,f^{\circ 2}(c_1),\,f^{\circ 3}(c_1)\}.}
\]
Equivalently, it is the boundary of a small regular neighbourhood of an arc
joining the two critical values \(f(c_1)\) and \(f(c_2)\), chosen disjoint from
the other three marked points.
\end{remark}

\begin{thebibliography}{9}

\bibitem{FKS}
T.~Firsova, J.~Kahn, and N.~Selinger,
\emph{On deformation spaces of quadratic rational maps},
arXiv:1606.00723.

\bibitem{HK}
E.~Hironaka and S.~Koch,
\emph{A disconnected deformation space of rational maps},
J. Mod. Dyn. \textbf{11} (2017), 409--423.

\end{thebibliography}

\end{document}
```